\section{Integrative Parameter Model of the Design Problem}
\label{sec:model}

Section~\ref{sec:method} described \emph{how} the human--AI--repository
loop is governed. This section describes \emph{what} it operates on:
the coupled parameter model of an electric VTOL UAV, the dimensioning
relations that implement it, and the constraint and verification checks
through which every iteration must pass. The two are complementary
contributions. The governance layer is what makes an AI agent safe to
employ; the integrative parameter model is what makes the agent
\emph{useful}, because it defines exactly which quantities a change
must be propagated to. An organization without heritage design tools
--- the situation of most emerging drone developers
(Section~\ref{sec:related}) --- needs both.

\subsection{Parameter taxonomy}
\label{sec:taxonomy}

The defining difficulty of electric VTOL conceptual design is that its
parameters do not separate cleanly into inputs and outputs. Wing area,
span, and airfoil are chosen by the designer, yet they are also
\emph{results} of the stall constraint at the converged weight; mission
endurance is a requirement, yet it is also what the closed battery mass
delivers. Treating such quantities as pure inputs is precisely the
error that produces internally inconsistent design cards --- as
happened in this project's own history (Section~\ref{sec:evolution}).
The model therefore classifies every quantity into one of five roles:

\begin{enumerate}[nosep,label=(\alph*)]
\item \textbf{Requirements} --- externally imposed and held fixed
  within an iteration: payload mass, mission segment durations, cruise
  speed, rate of climb, energy reserve, design ambient conditions.
\item \textbf{Free design parameters} --- the designer's (and the
  AI agent's) actual degrees of freedom: rotor diameter, blade count,
  pitch and solidity, aspect ratio, airfoil, thickness ratio,
  construction method and its structural knock-up factor, mass
  fractions, battery chemistry and cell count, component efficiencies.
\item \textbf{Coupled state variables} --- simultaneously input and
  output, resolved only by iteration: MTOW, battery mass, hover and
  cruise power, mission energy, wing area, span, and mean chord.
\item \textbf{Derived quantities} --- computed from the converged
  state, never configured: disk loading, induced and jet velocities,
  tip speed, dynamic pressures, lift-to-drag ratio, hover current,
  discharge C-rate, Joule heat, hinge moments, inertias, angular
  accelerations, pack temperature history.
\item \textbf{Constraints} --- inequalities the converged state must
  satisfy, each with an owner and a documented limit.
\end{enumerate}

Table~\ref{tab:params} lists the case-study instantiation. The
distinction is operationally important, not merely didactic: in the
repository, class (a) and (b) quantities live in commented
configuration files, class (c) and (d) quantities may only be produced
by the physics library, and class (e) quantities are asserted by the
test suite. The ``no magic numbers'' rule of
Section~\ref{sec:governance} is exactly the rule that keeps a class-(d)
quantity from being quietly hard-coded as if it were class-(a) --- the
failure mode behind the stall-speed inconsistency discussed later.

\begin{table}[htbp]
\centering
\caption{Parameter taxonomy of the case study. Requirements and free
parameters are configured with a written rationale; coupled and derived
quantities are computed only by the physics library; constraints are
asserted by the test suite.}
\label{tab:params}
\footnotesize
\begin{tabularx}{\textwidth}{@{}lXl@{}}
\toprule
\textbf{Role} & \textbf{Quantities (case study)} & \textbf{Lives in} \\
\midrule
Requirements &
$m_{\mathrm{pl}}=0.5$\,kg; $t_{\mathrm{hov}}=120$\,s,
$t_{\mathrm{tr}}=40$\,s, $t_{\mathrm{cr}}=900$\,s;
$V_{\mathrm{cr}}=20$\,m/s; RoC $=2.5$\,m/s;
$f_{\mathrm{res}}=1.20$; $\rho$, $T_\infty=40^{\circ}$C &
\code{config/} \\
\addlinespace
Free design parameters &
$D$, $N_b$, pitch, $\sigma$, FoM; AR, airfoil, $t/c$, $\lambda$,
$n_{\mathrm{ult}}$; construction method and $k$; $f_s,f_p,f_a,f_m$;
$e_{\mathrm{spec}}$, cell count $n_S$, $R_{\mathrm{pack}}$,
$f_{\mathrm{usable}}$; $\eta_p,\eta_m,\eta_{\mathrm{esc}}$ &
\code{config/} \\
\addlinespace
Coupled state variables &
$m_0$ (MTOW), $m_{\mathrm{batt}}$, $m_{\mathrm{wing}}$,
$P_{\mathrm{hov}}$, $P_{\mathrm{cr}}$, $E_{\mathrm{mis}}$,
$(W/S)^{\ast}$, $S$, $b$, $\bar{c}$ &
sizing loop \\
\addlinespace
Derived quantities &
$DL$, $v_i$, $V_{\mathrm{jet}}$, $q_{\mathrm{jet}}$,
$V_{\mathrm{tip}}$, $q_\infty$, $C_L$, $L/D$, $V_{\mathrm{MD}}$;
$I_{\mathrm{hov}}$, C-rate, $Q=I^2R_{\mathrm{pack}}$;
$M_{\mathrm{hinge}}$, $I_{xx}$, $\ddot\phi$, $\ddot\theta$;
$T_{\mathrm{batt}}(t)$, vent $\dot{m}$ &
\code{out/} hand-offs \\
\addlinespace
Constraints &
$V_{\mathrm{stall}}$, service ceiling, $T_{\mathrm{hov}}\le
T_{\max}$, C-rate $\le C_{\max}$, battery fraction, reserve energy,
$\ddot\theta\ge\ddot\theta_{\min}$, $T_{\mathrm{batt}}$ and
$T_{\mathrm{esc}}$ limits, packaging fit, CG range, $C_{D0}$ budget
(Table~\ref{tab:constraints}) &
test suite \\
\bottomrule
\end{tabularx}
\end{table}

\subsection{Parameter dependencies reported in the eVTOL literature}
\label{sec:dependencies}

Before mapping the couplings of the present vehicle it is worth
establishing which dependencies the eVTOL sizing literature already
treats as first-order, because the integrative model of
Section~\ref{sec:coupling} is a specialisation of them rather than a new
proposal. The consistent finding across eVTOL sizing studies is that
mass, geometry, aerodynamic performance, rotor loading, powertrain
choice, and transition feasibility feed back on one another rather than
being set independently, so sizing must be an iterative
multidisciplinary loop \citep{tyan2017,ugwueze2023,zhang2024evtol}. Two
consequences are specific to VTOL and recur throughout that literature.
First, the mission must be analysed at all critical points --- hover,
climb, transition, cruise, descent, and landing hover --- because
takeoff mass and energy capacity depend on every segment and not on
cruise alone \citep{tyan2017,ugwueze2022,hu2022}. Second, propulsion,
energy storage, and thermal management must be sized at component level,
since propeller, motor, ESC, and pack are strongly interdependent and
all scale with vehicle mass \citep{chahba2023,an2022,albamaestre2021}.

Table~\ref{tab:deps} collects the dependency structure in the
depends-on/drives form used by that literature, instantiated for a
single-propulsor tail-sitter. It is the qualitative statement of what
Fig.~\ref{fig:coupling} then makes concrete and
Section~\ref{sec:equations} quantifies.

\begin{table}[htbp]
\centering
\caption{Parameter dependencies for tail-sitter eVTOL conceptual design,
instantiated for a single-propulsor, battery-electric vehicle. The
depends-on/drives structure follows the eVTOL sizing literature cited in
each row; the parameter names are those of the present case study.}
\label{tab:deps}
\footnotesize
\begin{tabularx}{\textwidth}{@{}lXXX@{}}
\toprule
\textbf{Group} & \textbf{Primary parameters} & \textbf{Depends on} &
\textbf{Drives} \\
\midrule
Mission definition &
Payload, range/endurance, hover and transition budgets, cruise speed,
climb rate, ceiling, reserve &
Use case and top-level requirements \citep{he2025,lee2020} &
MTOW, stored energy, wing loading, rotor power
\citep{tyan2017,hu2022} \\
\addlinespace
Mass targets &
MTOW, structural/propulsion/avionics fractions, construction factor &
Mission energy, propulsion mass, airframe geometry, construction method
\citep{ugwueze2022,he2025} &
Required thrust, wing area, battery mass, transition feasibility
\citep{tyan2017,zhong2021} \\
\addlinespace
Wing and airframe geometry &
$W/S$, AR, span, area, chord, $C_{D0}$, $C_{L,\max}$, $L/D$ &
Cruise and stall constraints, transport/landing envelope, configuration
choice \citep{donateo2022,lee2020} &
Cruise power, stall margin, structural mass, transition corridor
\citep{hu2022,kubo2008} \\
\addlinespace
Rotor and propulsor &
Disk loading, diameter, blade count, solidity, pitch, thrust margin &
Hover power, takeoff weight, pack discharge capability
\citep{an2022,albamaestre2021,liu2024,qing2019} &
Hover efficiency, 1/rev forcing, jet dynamic pressure, control
authority, noise \citep{liu2024,govindarajan2018,sheng2015} \\
\addlinespace
Electric propulsion chain &
Motor and ESC rating, bus voltage, propeller--motor match &
Thrust and power required in every phase \citep{chahba2023} &
Propulsion mass, chain efficiency, thermal load, endurance
\citep{chahba2023} \\
\addlinespace
Battery system &
Specific energy and power, C-rate, DoD, pack resistance, voltage sag &
Mission profile, peak hover power, thermal limits
\citep{park2024,he2025} &
Battery mass, achievable hover/cruise power, cycle life, cooling need
\citep{park2024,zhao2024thermal} \\
\addlinespace
Thermal management &
Cooling path, airflow, vent and plate area, allowable temperature &
Heat generation, C-rate, hover duration
\citep{park2024,zhao2024thermal} &
Effective specific energy, thermal-system mass, degradation, safety
margin \citep{park2024,zhao2024thermal} \\
\addlinespace
Transition and control &
Pitch schedule, actuator and rate limits, vane and aileron authority,
time delay &
Airframe geometry, propulsion authority, wing aerodynamics
\citep{zhong2021,mcintosh2024,li2020} &
Feasible transition time, altitude loss, robustness, effector sizing
\citep{li2020,cong2024,banazadeh2016,marvakov2021} \\
\addlinespace
Volume and packaging integration &
Bay stack, pack volume, fuselage envelope, CG range &
Component envelopes, structural architecture, landing footprint
\citep{lee2020,govindarajan2018} &
Achievable pack size, static margin, structural mass
\citep{ugwueze2023} \\
\bottomrule
\end{tabularx}
\end{table}

Three findings from that literature are worth stating explicitly
because they bear on the constraint set of
Section~\ref{sec:checks}. Battery sizing constrained only by constant
specific energy is insufficient: voltage drop, temperature, depth of
discharge, and C-rate all belong in the conceptual loop
\citep{park2024}, and in the highest-power segments power density can
matter more than energy density \citep{thu2025}. Thermal management can
change feasibility outright rather than merely adding mass, since longer
hover demands substantially heavier cooling \citep{zhao2024thermal}.
And transition is a sizing constraint, not only a control problem:
imposing a minimum-transition-time constraint during sizing has been
shown to reduce takeoff weight by nearly 5\% relative to an unconstrained
initial design \citep{cong2024}, while the transition corridor widens
with available propulsive power and narrows with weight
\citep{tang2018,zhong2021}. For small battery-powered tail-sitters
specifically, higher battery mass fraction and lower wing loading
improve achievable performance \citep{wang2016}, and simultaneous
optimisation of wing loading, aspect ratio, and battery fraction against
energy consumption has been reported to raise the payload fraction
materially \citep{cinar2023}.

\subsection{The coupling map}
\label{sec:coupling}

Fig.~\ref{fig:coupling} is the integrative model: it shows which
factors influence which, and where the loops are. It follows the
depends-on/drives structure of Table~\ref{tab:deps}, and is drawn in the
sizing-flow form so that the primary path, the constraint and feedback
edges, and the dominant iterative loop are visually distinct. Two
circular dependencies dominate and must be resolved simultaneously.

The \textbf{closure loop} (heavy boxes, centre row) is the classical
electric-aircraft circularity: the battery is the heaviest single item,
its mass follows from mission energy, mission energy follows from
power, power follows from weight, and weight includes the battery.
Because the structural and subsystem masses are expressed as fractions
of MTOW, the loop admits an analytical step for weight given battery
mass, and the iteration reduces to a fixed point in
$m_{\mathrm{batt}}$ alone --- a contraction, because doubling the
vehicle weight does not double the required energy.

The \textbf{wing loop} enters through the stall constraint. The design
wing loading is bounded above by
$(W/S)_{\mathrm{stall}}=\tfrac{1}{2}\rho V_{\mathrm{stall}}^{2}
C_{L,\max}$, so wing area is set by weight, and weight is influenced by
wing mass. This is why the airfoil choice --- nominally a free
parameter --- propagates all the way into vehicle geometry: raising
$C_{L,\max}$ raises the admissible wing loading and \emph{shrinks} the
wing at constant weight. It is also why $C_{L,\max}$ must track the
airfoil selection in the same change; a stale value silently relaxes a
requirement rather than failing a test.

Two further couplings deserve note because they are easy to miss in a
sequential workflow, and both bit this project:

\begin{itemize}[nosep]
\item \textbf{Rotor diameter is a mass parameter, not only a
  propulsion parameter.} Disk loading is derived, $DL=m_0g/(\pi
  D^{2}/4)$, so a diameter change moves hover power, hence energy,
  hence battery mass, hence MTOW --- and simultaneously moves jet
  dynamic pressure, hence control-vane authority, hinge moments, servo
  torque requirement, and the 1/rev vibration forcing frequency.
\item \textbf{Pack internal resistance couples the thermal model to
  the mass closure.} $R_{\mathrm{pack}}$ sets both the Joule heat
  $I^{2}R$ that the battery bay must vent and the discharge efficiency
  $\eta_{\mathrm{bat}}$ that appears in the battery-mass relation.
  Treating the two independently permits a design that is thermally
  sized for a resistive pack while being mass-sized for an ideal one.
\end{itemize}

\begin{sidewaysfigure}
\centering
\hyphenpenalty=10000 \exhyphenpenalty=10000
\resizebox{\textheight}{!}{%
\begin{tikzpicture}[
  font=\scriptsize,
  bx/.style={draw, rounded corners=2pt, align=center, inner sep=1.4mm,
             minimum height=7mm, text width=30mm, fill=white},
  loop/.style={bx, very thick, fill=gray!15},
  a/.style={-{Stealth[length=1.8mm]}},
  fb/.style={-{Stealth[length=1.8mm]}, densely dashed}
]

% ================= band 0: external inputs =================
\node[bx, text width=34mm] (mis) at ( 4.4, 0.0)
  {\textbf{Mission requirements}\\payload, range / endurance,
   hover \& transition time, cruise speed, reserve};
\node[bx, text width=34mm] (ops) at (11.4, 0.0)
  {\textbf{Operational constraints}\\landing footprint, span,
   transport volume, turnaround};
\node[bx, text width=34mm] (arch) at (18.2, 0.0)
  {\textbf{Propulsion architecture}\\single ducted fan, jet vanes,
   battery-electric};
\node[bx, text width=32mm] (fail) at (24.6, 0.0)
  {\textbf{Failure cases}\\thrust loss, effector jam,
   required margins};

% ================= band 1: derived constraints / payload =================
\node[bx, text width=34mm] (dcon) at ( 3.4,-3.2)
  {\textbf{Derived constraints}\\stall speed, climb rate,
   ceiling, transition time};
\node[bx, text width=24mm] (pl) at ( 8.4,-3.2) {\textbf{Payload target}};

% ================= band 2: loadings, mass, subsystems =================
\node[bx] (wl) at ( 3.2,-6.2)
  {\textbf{Wing and power loading}};
\node[loop, text width=34mm] (m0) at ( 9.0,-6.2)
  {\textbf{Takeoff mass and mass breakdown}\\payload + structure +
   propulsion + energy + avionics};
\node[bx, text width=32mm] (str) at (14.6,-6.2)
  {\textbf{Structure and subsystems}\\members, landing loads,
   integration, avionics};
\node[bx, text width=32mm] (prop) at (19.8,-6.2)
  {\textbf{Propulsion sizing}\\installed power, motor / ESC,
   rotor diameter, disk loading};
\node[bx, text width=32mm] (ctl) at (25.0,-6.2)
  {\textbf{Control and stability}\\CG, vanes, ailerons,
   actuator limits};

% ================= band 3: performance chain =================
\node[bx] (wing) at ( 3.0,-10.2)
  {\textbf{Airframe sizing}\\area, span, aspect ratio, chord};
\node[bx] (aero) at ( 6.9,-10.2)
  {\textbf{Aerodynamics}\\$C_{L,\max}$, $C_{D0}$,
   $L/D$, stall margin};
\node[loop] (tr) at (10.8,-10.2)
  {\textbf{Transition feasibility}$^{\ast}$\\thrust margin,
   corridor width, altitude loss};
\node[bx] (cr) at (14.7,-10.2)
  {\textbf{Cruise performance}\\cruise drag,
   climb / cruise power};
\node[loop] (en) at (18.6,-10.2)
  {\textbf{Mission power and energy}\\hover + transition
   + cruise + reserve};
\node[bx] (lim) at (22.5,-10.2)
  {\textbf{Discharge and thermal limits}\\C-rate, DoD,
   pack resistance, temperature};
\node[bx, text width=28mm] (vol) at (26.2,-10.2)
  {\textbf{Volume integration}\\pack volume, bay stack,
   fuselage envelope};

% ================= band 4/5: energy storage =================
\node[loop, text width=34mm] (es) at (18.6,-14.0)
  {\textbf{Energy storage sizing}\\energy capacity and
   power capability};
\node[bx, text width=34mm] (esm) at (18.6,-16.6)
  {\textbf{Energy storage mass}};

% ================= primary sizing flow =================
\draw[a] (mis.south)  -- (dcon.north);
\draw[a] (mis.south)  -- (pl.north);
\draw[a] (arch.south) -- (prop.north);
\draw[a] (dcon.south) -- (wl.north);
\draw[a] (pl.south)   -- (m0.north);
\draw[a] (m0)   -- (str);
\draw[a] (str)  -- (prop);
\draw[a] (prop) -- (ctl);
\draw[a] (wl.south)   -- (wing.north);
\draw[a] (wing) -- (aero);
\draw[a] (aero) -- (tr);
\draw[a] (tr)   -- (cr);
\draw[a] (cr)   -- (en);
\draw[a] (en)   -- (lim);
\draw[a] (prop.south) -- (en.north);
\draw[a] (en.south)   -- (es.north);
\draw[a] (lim.south)  -- (es.east);
\draw[a] (es.south)   -- (esm.north);
% wing geometry informs the structural member model
\draw[a] (aero.north) -- (6.9,-8.2) -- (14.6,-8.2) -- (str.south);

% ================= constraints and feedback =================
\draw[fb] (ops.south)  -- (pl.north);
\draw[fb] (ops.south)  -- (m0.north);
\draw[fb] (fail.south) -- (ctl.north);
\draw[fb] (ctl.south)  -- (vol.north);
\draw[fb] (vol.south)  -- (es.east);
% converged mass drives the loadings back
\draw[fb] (m0.west) -- (wl.east);
% transition feasibility constrains the mass breakdown
\draw[fb] (tr.north) -- (10.8,-8.6) -- (9.0,-8.6) -- (m0.south);
% the dominant loop closes: energy storage mass re-enters the mass breakdown
\draw[fb] (esm.west) -- (0.4,-16.6) -- (0.4,-4.9) -- (9.0,-4.9) -- (m0.north);

% ================= legend =================
\begin{scope}
  \node[draw, rounded corners=2pt, minimum width=64mm, minimum height=28mm,
        anchor=north west] at (1.4,-13.4) {};
  \draw[a] (1.8,-14.0) -- (3.1,-14.0);
  \node[anchor=west] at (3.3,-14.0) {solid: primary sizing flow};
  \draw[fb] (1.8,-14.7) -- (3.1,-14.7);
  \node[anchor=west] at (3.3,-14.7) {dashed: constraint / feedback};
  \node[loop, text width=9mm, minimum height=4mm, inner sep=0.6mm]
        at (2.45,-15.4) {};
  \node[anchor=west] at (3.3,-15.4) {thick grey: dominant iterative loop};
  \node[anchor=north west, align=left] at (1.7,-15.7)
    {$^{\ast}$ general model; the case study budgets transition
     energetically\\ but does not size its corridor
     (Section~\ref{sec:discussion})};
\end{scope}
\end{tikzpicture}%
}
\caption{Integrative sizing and coupling map for tail-sitter eVTOL
conceptual design, in the depends-on/drives structure of
Table~\ref{tab:deps} and instantiated for the case-study vehicle
(diagram topology after \c{C}.~U\c{c}ler). Solid edges are the primary
sizing flow; dashed edges are constraints and feedback; the four heavy
grey boxes are the dominant iterative loop, which closes when the energy
storage mass re-enters the mass breakdown. Representative edges are
drawn; the complete dependency set is enforced in code.}
\label{fig:coupling}
\end{sidewaysfigure}

\subsection{Dimensioning relations}
\label{sec:equations}

The relations below are the conceptual-level physics the pipeline
implements. They are deliberately first-order and transparent: the
methodology's claim is not fidelity but traceability, and every symbol
resolves either to a configured parameter with a written rationale or
to another equation in this set. They are used in the form collected by
\citet{poelma2024}, which consolidates the Gudmundsson/Sadraey
constraint-diagram relations, the Raymer weight regressions, and the
VTOL power-to-weight formulation of \citet{tyan2017} into one
conceptual-sizing procedure for electric VTOL UAVs; the section
references the primary sources alongside.

\paragraph{Wing loading and thrust loading.} With a parabolic drag
polar $C_D=C_{D0}+kC_L^{2}$, $k=1/(\pi\,\mathrm{AR}\,e)$, and
$q=\tfrac{1}{2}\rho V^{2}$, each flight phase gives a required thrust
loading as a function of wing loading \citep{gudmundsson2014,sadraey2013}:
\begin{align}
\text{cruise:}\quad
  \frac{T}{W} &= \frac{q\,C_{D0}}{W/S} + \frac{k\,(W/S)}{q},
  \qquad q=\tfrac{1}{2}\rho V_{\mathrm{cr}}^{2},
  \label{eq:smd}\\
\text{climb:}\quad
  \frac{T}{W} &= \frac{\mathrm{RoC}}{V} + \frac{q\,C_{D0}}{W/S}
                 + \frac{k\,(W/S)}{q},
  \qquad V=\sqrt{\frac{2}{\rho}\sqrt{\frac{k}{3C_{D0}}}\,\frac{W}{S}},
  \label{eq:climb}\\
\text{ceiling:}\quad
  \frac{T}{W} &= \frac{\mathrm{RoC}_{\mathrm{ceil}}}{V_{\mathrm{mp}}}
                 + 4\sqrt{\frac{k\,C_{D0}}{3}},
  \label{eq:ceiling}\\
\text{max speed:}\quad
  \frac{T}{W} &= \frac{\rho V_{\max}^{2}C_{D0}}{2\,(W/S)}
                 + \frac{2k\,(W/S)}{\rho V_{\max}^{2}} .
  \label{eq:vmax}
\end{align}
The stall requirement is a vertical boundary rather than a curve,
\begin{equation}
\left(\frac{W}{S}\right)_{\mathrm{stall}}
  = \tfrac{1}{2}\,\rho\,V_{\mathrm{stall}}^{2}\,C_{L,\max},
\label{eq:stall}
\end{equation}
and the design point minimises the upper envelope of
Eqs.~(\ref{eq:smd})--(\ref{eq:vmax}) inside the feasible region,
\begin{equation}
\begin{aligned}
(W/S)^{\ast}
  &= \operatorname*{arg\,min}_{\,W/S\,\le\,(W/S)_{\mathrm{stall}}}\;
     \max_{\text{phases}} \frac{T}{W}\!\left(\frac{W}{S}\right),\\[2pt]
(T/W)^{\ast}
  &= \max_{\text{phases}}\frac{T}{W}\Big((W/S)^{\ast}\Big),
\end{aligned}
\label{eq:designpoint}
\end{equation}
which simultaneously minimises propulsion size and battery mass. Wing
geometry then follows from the converged weight,
\begin{equation}
S = \frac{m_0\,g}{(W/S)^{\ast}},
\qquad
b = \sqrt{\mathrm{AR}\cdot S},
\qquad
\bar{c} = \frac{S}{b},
\label{eq:winggeom}
\end{equation}
with wing structural mass from Raymer's general-aviation regression
\citep{raymer2018}, scaled by a material factor
$k_{\mathrm{mat}}$ ($0.70$ for CFRP):
\begin{equation}
m_{\mathrm{wing}} = k_{\mathrm{mat}}\;0.036\,
  S^{0.758}\,
  \left(\frac{\mathrm{AR}}{\cos^{2}\Lambda}\right)^{0.6}
  q^{0.006}\,\lambda^{0.04}
  \left(\frac{100\,t/c}{\cos\Lambda}\right)^{-0.3}
  \left(n_{\mathrm{ult}}W_0\right)^{0.49},
\label{eq:raymer}
\end{equation}
evaluated in the imperial units of the original regression and
converted back to SI.

\paragraph{Hover and vertical flight.} For a tail-sitter the single
rotor carries the entire weight, so disk loading is derived rather than
chosen:
\begin{equation}
DL = \frac{m_0\,g}{\pi D^{2}/4}.
\label{eq:dl}
\end{equation}
Momentum (actuator-disk) theory with a figure of merit gives the hover
shaft power, and the vertical-climb case adds blade profile drag and
the parasite drag of fuselage and wing in the vertical airstream:
\begin{align}
\left(\frac{P}{W}\right)_{\mathrm{hov}}
  &= \frac{1}{\mathrm{FoM}}\sqrt{\frac{DL}{2\rho}},
  \label{eq:hover}\\
\left(\frac{P}{W}\right)_{\mathrm{climb}}
  &= \underbrace{\frac{\mathrm{RoC}}{2}
     + \frac{1}{2}\sqrt{\mathrm{RoC}^{2}+\frac{2\,DL}{\rho}}}
     _{\text{induced}+\text{climb}}
   + \underbrace{\frac{\rho\,V_{\mathrm{tip}}^{3}\,\sigma\,C_{d}}{8\,DL}}
     _{\text{blade profile}}
   \notag \\
  &\qquad
   + \underbrace{\frac{\rho\,\mathrm{RoC}^{3}}{DL}}_{\text{fuselage}}
   + \underbrace{\frac{\rho\,\mathrm{RoC}^{3}}{S_{r}\,(W/S)}}_{\text{wing}},
  \label{eq:vclimb}
\end{align}
with tip speed from an empirical rotor speed--diameter relation,
$N = k_{\mathrm{rpm}}D^{\,n}$ and
$V_{\mathrm{tip}} = \pi N D/60$. Electrical power adds the drive chain
and the continuous avionics (hotel) load,
\begin{equation}
P_{\mathrm{hov}}
  = \frac{(P/W)_{\mathrm{vtol}}\;m_0\,g}
         {\eta_m\,\eta_{\mathrm{esc}}} + P_{\mathrm{hotel}},
\qquad
(P/W)_{\mathrm{vtol}} = \max\!\left[
  \left(\tfrac{P}{W}\right)_{\mathrm{hov}},
  \left(\tfrac{P}{W}\right)_{\mathrm{climb}}\right],
\label{eq:phov}
\end{equation}
while wing-borne cruise power uses the lift-to-drag model,
\begin{equation}
P_{\mathrm{cr}} = \frac{m_0\,g\,V_{\mathrm{cr}}}{(L/D)\;\eta_{\mathrm{tot}}}
                  + P_{\mathrm{hotel}},
\qquad
\eta_{\mathrm{tot}} = \eta_p\,\eta_m\,\eta_{\mathrm{esc}}.
\label{eq:pcr}
\end{equation}

\paragraph{Energy, battery, and mass closure.} Transitions are billed
conservatively at hover power, and the reserve is a multiplier on the
total:
\begin{equation}
E_{\mathrm{mis}} = \frac{f_{\mathrm{res}}}{3600}\Big[
  P_{\mathrm{hov}}\,(t_{\mathrm{hov}}+t_{\mathrm{tr}})
  + P_{\mathrm{cr}}\,t_{\mathrm{cr}}\Big].
\label{eq:energy}
\end{equation}
Battery mass follows from pack-level specific energy, discharge
efficiency, and usable depth of discharge \citep{tyan2017}, and MTOW
closes analytically against the mass fractions:
\begin{equation}
m_{\mathrm{batt}} = \frac{E_{\mathrm{mis}}}
  {e_{\mathrm{spec}}\;\eta_{\mathrm{bat}}\;f_{\mathrm{usable}}},
\qquad
m_0 = \frac{m_{\mathrm{pl}} + m_{\mathrm{batt}}}
           {1 - (f_s + f_p + f_a + f_m)} .
\label{eq:closure}
\end{equation}
Equations~(\ref{eq:dl})--(\ref{eq:closure}) are the closure loop of
Fig.~\ref{fig:coupling}, iterated to
$|m_{\mathrm{batt}}^{(i+1)}-m_{\mathrm{batt}}^{(i)}| < \varepsilon$.
The structural fraction is scaled by the construction method,
$f_s = k\,f_{s,\mathrm{base}}$, which makes $k$ the single most
sensitive mass parameter in the model. The converged operating point
gives the peak discharge rate and the pack current used by the
electrical and thermal stages:
\begin{equation}
C_{\mathrm{peak}}
  = \frac{\max(P_{\mathrm{hov}},P_{\mathrm{cr}})}
         {m_{\mathrm{batt}}\,e_{\mathrm{spec}}},
\qquad
I_{\mathrm{hov}} = \frac{P_{\mathrm{hov}}}{n_S\,V_{\mathrm{cell}}} .
\label{eq:crate}
\end{equation}

\paragraph{Control authority.} Jet-vane authority is sized from hover
thrust. With $T=m_0g$ the induced velocity, fully developed jet
velocity, and vane dynamic pressure are
\begin{equation}
v_i = \sqrt{\frac{T}{2\rho A}},
\qquad V_{\mathrm{jet}} = 2v_i,
\qquad q_{\mathrm{jet}} = \tfrac{1}{2}\rho V_{\mathrm{jet}}^{2}
  = \frac{T}{A},
\label{eq:jet}
\end{equation}
so vane force, control moment, angular acceleration, and the servo
torque requirement follow as
\begin{equation}
F_v = q_{\mathrm{jet}} S_v C_{l\alpha}\delta,
\qquad
\ddot\theta = \frac{n_v F_v \ell}{I_{\theta}},
\qquad
M_{\mathrm{hinge}} = q_{\mathrm{jet}} S_v c_v C_{h},
\qquad
M_{\mathrm{servo}} = M_{\mathrm{hinge}}\,\mathrm{SF}\,f_m .
\label{eq:vane}
\end{equation}
Equation~(\ref{eq:jet}) exposes why a second roll effector is
structurally necessary rather than merely prudent: $q_{\mathrm{jet}} =
T/A$ is \emph{linear} in thrust, and cruise thrust is smaller than
hover thrust by roughly the factor $L/D$, so vane authority collapses
in cruise. Ailerons instead use free-stream dynamic pressure, with
Glauert flap effectiveness
$\tau = 1-(\theta_f-\sin\theta_f)/\pi$,
$\theta_f = \arccos(2c_f/c-1)$:
\begin{equation}
\ddot\phi = \frac{2\,q_\infty\,S_a\,(C_{L\alpha}\tau)\,\delta\,y_a}
                 {I_{xx}} .
\label{eq:aileron}
\end{equation}
The two effectors are therefore complementary by construction: each is
strongest exactly where the other is weakest, and both are checked
against the same authority floor $\ddot\theta_{\min}$.

\paragraph{Thermal paths.} The two hover heat loads are derived from
the converged operating point, not configured:
$Q_{\mathrm{esc}} = P_{\mathrm{hov}}(1-\eta_{\mathrm{esc}})$ and
$Q_{\mathrm{batt}} = I_{\mathrm{hov}}^{2}R_{\mathrm{pack}}$. Forced
convection on the inflow-washed cold plate and the vented battery bay
give the steady sizing relations
\begin{equation}
T_{\mathrm{plate}}(A) = T_\infty + \frac{Q_{\mathrm{esc}}}{C_h A^{3/4}},
\qquad
\Delta T_{\mathrm{air}} = \frac{Q_{\mathrm{batt}}}{\dot{m}\,c_p},
\qquad
\dot{m} = \rho\,V_{\mathrm{vent}}\,A_{\mathrm{vent}} .
\label{eq:thermalss}
\end{equation}
Because the vent check is air-side only, it cannot see the pack's own
temperature rise; a lumped-capacitance mission transient with
$C_{\mathrm{th}} = m_{\mathrm{batt}}c_{p,\mathrm{pack}}$ and
$\tau = C_{\mathrm{th}}/(UA)$ is therefore integrated over the two
vertical legs and the cruise segment between them,
\begin{equation}
\Delta T_{\mathrm{leg}} = \frac{Q_{\mathrm{hov}}\,t_{\mathrm{leg}}}
                               {C_{\mathrm{th}}},
\qquad
T(t) = T_{\mathrm{ss}} + \big(T_0-T_{\mathrm{ss}}\big)e^{-t/\tau},
\qquad
T_{\mathrm{ss}} = T_\infty + \frac{Q_{\mathrm{cr}}}{UA},
\label{eq:transient}
\end{equation}
and evaluated at a bracket of pack resistances. This is the check that
exposed the design's second standing finding
(Section~\ref{sec:design-results}).

\subsection{Constraints and the iteration check flow}
\label{sec:checks}

Table~\ref{tab:constraints} lists the constraint set with its source
and enforcement point. Two properties matter methodologically. First,
each limit is \emph{configured with a rationale} while the quantity
compared against it is \emph{derived} --- so a design change cannot
relax a requirement by accident, only by an explicit, reviewable edit
to a configuration file. Second, each constraint has a named
enforcement mechanism, ranging from a hard exception that aborts the
sizing loop to a reported margin that is allowed to be negative
provided it is carried as a standing finding.

\begin{table}[htbp]
\centering
\caption{Constraint set of the case study, with the enforcement
mechanism for each. ``Hard'' aborts the pipeline; ``pinned'' is
asserted by the regression suite; ``reported'' is a margin that may be
negative only as a documented standing finding.}
\label{tab:constraints}
\footnotesize
\begin{tabularx}{\textwidth}{@{}Xllc@{}}
\toprule
\textbf{Constraint} & \textbf{Limit (case study)} &
\textbf{Owner} & \textbf{Type} \\
\midrule
Battery fraction (closure divergence guard) &
$m_{\mathrm{batt}}/m_0 \le 0.60$ & sizing loop & hard \\
Hover thrust within COTS rotor capability &
$T_{\mathrm{hov}} \le T_{\max}=35$\,N & sizing study & hard \\
Mass-fraction sum consistency &
$\sum f_i + f_{\mathrm{batt}} + f_{\mathrm{pl}} = 1$ &
sizing loop & pinned \\
Stall speed &
$V_{\mathrm{stall}} \le 12$\,m/s via Eq.~(\ref{eq:stall}) &
wing stage & pinned \\
Service ceiling &
$h \ge 3000$\,m at $\mathrm{RoC}=0.5$\,m/s &
size matching & pinned \\
Peak discharge rate &
$C_{\mathrm{peak}} \le 25$\,C & sizing loop & pinned \\
Mission energy reserve &
$f_{\mathrm{res}} = 1.20$ & sizing loop & pinned \\
Angular acceleration, all axes, hover and cruise &
$\ddot\theta \ge 30^{\circ}$/s$^{2}$ & vane, aileron stages & pinned \\
Vibration isolation window at 1/rev &
transmissibility target at $\sim$204\,Hz & isolation stage & pinned \\
ESC junction temperature &
$T_{\mathrm{esc}}$ limit at $T_\infty=40^{\circ}$C &
thermal stage & reported \\
Battery pack temperature (steady and transient) &
$T_{\mathrm{batt}} \le 60^{\circ}$C & thermal stage & reported \\
Packaging fit, bay stack, CG range &
rigid-part axial lengths, static margin &
fuselage stages & pinned \\
Cruise drag budget &
$C_{D0} \le 0.025$ (build-up vs.\ configured) &
fuselage stage & reported \\
Mass allocation per component group &
as-selected mass $\le$ fraction allocation &
freeze stage & reported \\
Cross-consumer reference consistency &
CFD $l_{\mathrm{ref}}$, $A_{\mathrm{ref}}$, CofR vs.\ hand-off &
test suite & pinned \\
\bottomrule
\end{tabularx}
\end{table}

Fig.~\ref{fig:checkflow} traces how these checks are actually used
within one iteration. The flow has three distinct failure exits, and
keeping them distinct is what makes the framework usable as a review
instrument. A \emph{hard} failure means the concept does not close at
all and the sizing loop refuses to produce a number --- the response is
to change a free parameter or, where the assumption itself was wrong,
to amend it by decision record. A \emph{pinned} failure means the
design moved: if the move was intended, the pins, downstream consumer
references, decision record, and changelog are updated in the same
commit, so the diff documents the new design point; if it was not, the
merge is blocked because the design has drifted silently. A
\emph{reported} negative margin does not block anything, but it cannot
be silenced either: it is emitted by the final pipeline stage on every
run until it is retired.

\begin{figure}[htbp]
\centering
\resizebox{0.90\textwidth}{!}{%
\begin{tikzpicture}[
  font=\footnotesize,
  proc/.style={draw, rounded corners=2pt, align=center, inner sep=2mm,
               minimum height=9mm, fill=gray!8},
  dec/.style={draw, diamond, aspect=2.4, align=center, inner sep=0.6mm,
              fill=orange!15},
  bad/.style={draw, rounded corners=2pt, align=center, inner sep=2mm,
              fill=red!12, minimum height=9mm},
  good/.style={draw, rounded corners=2pt, align=center, inner sep=2mm,
               fill=green!14, minimum height=9mm},
  a/.style={-{Stealth[length=2mm]}},
  fb/.style={-{Stealth[length=2mm]}, densely dashed, thick},
  el/.style={font=\scriptsize, inner sep=1pt}
]
% ---- spine ----
\node[proc] (cfg)   at (4.6,  0.0) {\code{config/*.yaml}: requirements,\\free parameters, margins};
\node[proc] (smd)   at (4.6, -1.7) {Size-matching diagram\\$(W/S)^{\ast}$, $(T/W)^{\ast}$};
\node[proc] (loop)  at (4.6, -3.4) {Mass-closure iteration\\Eq.~(\ref{eq:dl})--(\ref{eq:closure})};
\node[dec]  (conv)  at (4.6, -5.3) {converged?};
\node[dec]  (guard) at (4.6, -7.3) {guards pass?};
\node[proc] (disc)  at (4.6, -9.2) {Discipline stages: authority,\\isolation, thermal, packaging};
\node[proc] (frz)   at (4.6,-11.1) {Component freeze +\\as-selected re-solve};
\node[proc] (marg)  at (4.6,-12.9) {Margins table +\\standing findings};
\node[dec]  (pins)  at (4.6,-14.9) {pins match?};
\draw[a] (cfg)  -- (smd);
\draw[a] (smd)  -- (loop);
\draw[a] (loop) -- (conv);
\draw[a] (conv) -- node[el,right]{yes} (guard);
\draw[a] (guard)-- node[el,right]{yes} (disc);
\draw[a] (disc) -- (frz);
\draw[a] (frz)  -- (marg);
\draw[a] (marg) -- (pins);

% ---- right-hand branches ----
\node[dec]  (div)   at (11.0, -5.3) {$m_{\mathrm{batt}}/m_0>0.6$?};
\node[bad]  (hard)  at (16.6, -6.4) {\textbf{hard failure}\\the concept does not\\close: no number is\\produced};
\node[dec]  (intent)at (11.0,-14.9) {intended?};
\node[proc] (upd)   at (11.0,-17.9) {update pins, consumer refs,\\ADR, changelog\\in the same commit};
\node[bad]  (drift) at (16.6,-14.9) {\textbf{silent drift}\\merge blocked};
\node[good] (merge) at (16.6,-17.9) {CI green $\rightarrow$\\human review $\rightarrow$\\merge and tag};

\draw[a] (conv.east)  -- node[el,above]{no} (div.west);
\draw[a] (div.north)  -- node[el,right]{no} (11.0,-2.5)
                      -- (7.9,-2.5) -- (7.9,-3.4) -- (loop.east);
\draw[a] (div.east)   -- node[el,above]{yes} (hard.west);
\draw[a] (guard.east) -- node[el,above]{no} ++(1.3,0) |- (hard.west);
\draw[a] (pins.east)  -- node[el,above]{no} (intent.west);
\draw[a] (intent.east)-- node[el,above]{no} (drift.west);
\draw[a] (intent.south)-- node[el,right]{yes} (upd.north);
\draw[a] (upd.east)   -- (merge.west);
\draw[a] (pins.south) -- node[el,right]{yes} ++(0,-1.0) -| (merge.north);

% ---- design-revision feedback channel (over the top, back to config) ----
\draw[fb] (hard.north) -- (16.6,1.8) -- (4.6,1.8) -- (cfg.north);
\node[font=\scriptsize, above] at (10.6,1.85)
  {revise a free parameter, or amend the assumption by ADR, and re-converge};

\node[align=left, font=\scriptsize] at (11.4,-12.9)
  {reported margins may be negative \emph{only}\\
   as standing findings: re-emitted by the\\
   final stage on every run};
\draw[fb] (marg.east) -- ++(0.9,0);
\end{tikzpicture}%
}
\caption{How the constraint checks of Table~\ref{tab:constraints} are
used inside one design iteration. Three failure exits are kept
deliberately distinct: a hard guard failure means the concept does not
close, an unexpected regression-pin failure means the design drifted
silently and blocks the merge, and a negative reported margin is
admissible only as a documented standing finding.}
\label{fig:checkflow}
\end{figure}

\subsection{Running the model in the other direction}
\label{sec:duality}

A practical consequence of the taxonomy is that the same coupling graph
serves several distinct questions, depending on which quantities are
declared fixed. In the case study the mission is a requirement and the
battery closes to it, so endurance is an input and pack mass is an
output. Reversing that assignment --- fixing a procured pack and
solving for achievable hover and cruise time --- traverses
Eqs.~(\ref{eq:energy})--(\ref{eq:closure}) in the opposite direction
and turns endurance into the output, which is the form needed once
hardware is on the bench. Likewise, wing area is an output of the
stall-limited closure here, but a fixed-span requirement (transport
container, launcher rail) makes span the input and stall speed the
output. The model does not change; only the partition between classes
(a) and (c) does.

This matters for the methodology because the partition is declared in
configuration rather than hard-coded in the equations, so switching the
question is a configuration change that the same verification harness
still gates. It is also the mechanism by which the framework absorbs
procurement reality: the component-freeze stage
(Section~\ref{sec:iteration}) is precisely a re-partition, in which
quantities that were free parameters during sizing become fixed inputs
taken from a datasheet.
